\outclassframe{
\begin{frame}{\ }\relax
    Very few people make full-fledged charts in tikZ. But it is much more common to add tikZ as text decoration element.
\end{frame}
}


\begin{frame}[fragile]{Introduction}\relax

TikZ is often used not as ``independent picture'', but as a part of the presentation or document. 

\end{frame}

\begin{frame}[fragile]{Example: CV\magicPage}\relax
    \begin{columns}
        \begin{column}{0.45\textwidth}
              \only<1>{
              \expandafter{\inputminted[firstline= 16, lastline=20]{latex}{sec01/code/cv_examples.tex}
                }
                }
                \only<2>{\expandafter{\inputminted[firstline= 22, lastline=32]{latex}{sec01/code/cv_examples.tex}}}
        \end{column}
        \begin{column}{0.45\textwidth}
             \includegraphics[width=\textwidth, keepaspectratio]{cv_examples}
        \end{column}
    \end{columns}

    % \cprotect\twocolImg{
    % \only<1>{\inputminted[firstline= 16, lastline=20]{latex}{sec01/code/cv_examples.tex}}
    % \only<2>{\inputminted[firstline= 22, lastline=32]{latex}{sec01/code/cv_examples.tex}}
    % }{cv_examples}  
    
\end{frame}

\begin{frame}[fragile]{``Magic''}\relax

    How to produce ``magic'' \only<1>{\tikz[overlay, remember picture] \draw[->] (0, 0.1) to[out=0,in=180] (endM);}

\begin{tikzpicture}[remember picture,overlay,shift={(current page.north east)}]
    \node[anchor=north east,xshift=-0cm,yshift=-0cm](endM) {%
        {\includegraphics[width=1cm]{black_magic}}%
};
\end{tikzpicture}%

\inpause
\begin{minted}{latex}
\begin{tikzpicture}[remember picture,overlay,shift={(current page.north east)}]
\node[anchor=north east,xshift=-0cm,yshift=-0cm](endM) {%
    {\includegraphics[width=1cm]{images/magic}}%
};
\end{tikzpicture}
\end{minted}

\inpause
Notice \ccol{shift=\{(current page.north east)\}}
\end{frame}



\begin{frame}[fragile]{arrow to magic}\relax

How to produce this \tikz[overlay, remember picture] \node (arrstart) {}; arrow \tikz[overlay, remember picture] \node (startMMM) {};

\begin{tikzpicture}[remember picture,overlay,shift={(current page.north east)}]
    \node[anchor=north east,xshift=-0cm,yshift=-0cm](endMMM) {%
        {\includegraphics[width=1cm]{black_magic}}%
};
\end{tikzpicture}%

\tikz[overlay, remember picture] \draw[->, ultra thick] ($(startMMM)+(0,0.1)$) to[out=0,in=180] (endMMM);
\only<1>{\begin{tikzpicture}[overlay, remember picture]
     \path[->, ultra thick, name path=point c] ($(startMMM)+(0,0.1)$) to[out=0,in=180] (endMMM);
     
     \path[->, name path=point a] ($(startMMM)+(5,-1)$) -- +(-1, 3);
     \path[name intersections={of=point c and point a, by={intersection}}];
     \draw[thick, thick, ->] ($(arrstart)-(0.4ex,-0.6ex)$) to[out=45, in=110] (intersection);
\end{tikzpicture}}

\inpause

\begin{minted}{latex}
% remember position of (startM) node (and (endM) node from previous slide)
How to produce this arrow \tikz[overlay, remember picture] \node (startM) {};

% go from (startM) to (endM)
\tikz[overlay, remember picture] \draw[->, ultra thick] ($(startM)+(0,0.1)$) to[out=0,in=180] (endM);

\end{minted}

\skfootnote{What about arrow pointed on the thick arrow? It was produced with intersections library}
     
\end{frame}


\begin{frame}[fragile]{``Common belief''\magicPage}\relax
     \begin{center}
        \begin{tikzpicture}
             \node[align=center] (0,0) {
             \huge \LaTeX\ is only for use\\ \huge  in academic area
             };
             \uncover<2,3>{\node[rotate=30, bottom color=red!50, top color=red!50] (0,0) {\Huge WRONG};}
        \end{tikzpicture}
         
    \end{center}
    
\pause\pause was produce by 
\begin{minted}{latex}
\begin{tikzpicture}
    \node[align=center] (0,0) {
    \huge \LaTeX\ is only for use\\ \huge  in academic area
    };
    \uncover<2,3>{\node[rotate=30, bottom color=red!50, top color=red!50] (0,0) {\Huge WRONG}};
\end{tikzpicture}
\end{minted}
    
\end{frame}


\cprotect\inclassframe{
\begin{frame}[fragile]{\exFrame{Try to use other people's work}}{lvl 1}\relax

    \begin{enumerate}
        \item You have the code as below. Copy it to you document and make sure it is compiling. ``Play'' with the params and of \ccol\pictofraction\ and as the result give the enumerate list in pdf of what all params means. Ex: ``1. printed symbol, 2...''
         \item Go to \url{https://tikz.dev/library-mindmaps} and try to create your own mindmap
    \end{enumerate}

    \inputminted[firstline=7, lastline=17]{latex}{sec01/code/exExtLvl1.tex}
    \inputminted[firstline=21, lastline=22]{latex}{sec01/code/exExtLvl1.tex}
%7-17, 21-22
    %  \includegraphics[width=0.4\textwidth]{exLvl1}
\end{frame}

\begin{frame}[fragile, allowframebreaks]{\exFrame{Try to use other people's work}}{lvl 2}\relax
You copy a code, that produce a ``day cyrcle'' (from AltaCV). But some libraries and packages are missing. See logs and google errors, make this code works :)

\includegraphics[width=0.8\textwidth]{exExtLvl2full}

     \inputminted[fontsize=\tiny]{latex}{sec01/code/exExtLvl2notfull.tex}
    %  \includegraphics[width=0.4\textwidth]{exLvl2}
\end{frame}

}



